\section{Introduction}

Speech is increasingly converted into operational signals. Contact-center and anti-fraud workflows use ASR hypotheses as inputs to routing, escalation, compliance review, and case-handling decisions. In these settings, transcript accuracy is necessary but not sufficient: a small ambiguity around a decision-critical detail can change whether a case is escalated, routed to routine review, or held for conservative manual assessment.

This paper focuses on high-stakes speech-driven decision systems where errors affect spoken decision atoms: negation, amount, actor, action, time, intent, uncertainty, and scam pattern. A transcript can remain close under WER or CER while a plausible ASR alternative changes the downstream action. The central evaluation question is therefore not only how similar the transcript is to a reference, but whether plausible transcript alternatives would preserve the decision.

We introduce Counterfactual Decision-Stability ASR (CDS-ASR), a downstream-aware ASR evaluation framework. CDS-ASR identifies risk atoms, constructs bounded plausible ASR counterfactual variants, scores decision instability with CEIS, and evaluates conservative recovery policies through aggregate replay. The motivating case is anti-fraud triage, but the method is framed as speech evaluation: it tests whether ASR hypotheses remain stable with respect to downstream actions.

The paper makes four contributions. First, it formulates decision-stability evaluation for ASR in high-stakes speech-to-decision workflows. Second, it defines risk atoms as decision-critical transcript spans and uses plausible ASR alternatives to stress the downstream decision function. Third, it evaluates CEIS against WER, CER, and SRES using human-reviewed decision-change labels from 30 reviewed rows and 90 clustered model assessments. Fourth, it reports aggregate policy replay as a scoped intervention analysis, with thresholds treated as diagnostic operating points rather than deployment rules.

All empirical claims are tied to aggregate artifacts. Raw audio, raw transcripts, selected row identifiers, reviewer notes, and transcript-bearing runtime logs remain outside the public release boundary.
